% Generated draft: requires mathematical and source review.
\documentclass[11pt]{article}
\usepackage[margin=25mm]{geometry}
\usepackage{fontspec}
\setmainfont{Latin Modern Roman}
\usepackage{amsmath,amssymb,mathtools,unicode-math}
\setmathfont{Latin Modern Math}
\usepackage{enumitem,needspace,xurl,hyperref}
\setlength{\parindent}{0pt}
\setlength{\parskip}{5pt}
\setlength{\emergencystretch}{3em}
\newcommand{\N}{\mathbb{N}}
\newcommand{\Z}{\mathbb{Z}}
\newcommand{\Q}{\mathbb{Q}}
\newcommand{\R}{\mathbb{R}}
\newcommand{\C}{\mathbb{C}}
\newcommand{\F}{\mathbb{F}}
\newcommand{\A}{\mathbb{A}}
\newcommand{\PP}{\mathbb{P}}
\newcommand{\E}{\mathbb{E}}
\newcommand{\eps}{\varepsilon}
\newcommand{\dd}{\,\mathrm{d}}
\begin{document}
\section*{UnsolvedMath ID: 30003469}
\textbf{Source problem code:} catalogue entry
\subsection{Definitions and mathematical statement}
Let \(d,k\in\Z_{\geq 1}\), and write
\[
[0,k]^d=\{x\in\R^d:0\leq x_i\leq k\text{ for }1\leq i\leq d\}.
\]
A lattice polytope is the convex hull of a nonempty finite subset of \(\Z^d\). Its graph has the vertices of the polytope as nodes and its one-dimensional faces as edges. Its combinatorial diameter \(\operatorname{diam}_{\mathrm{gr}}(P)\) is the maximum, over pairs of vertices, of the minimum number of edges in a path joining them; a singleton has diameter zero.

Under the combinatorial interpretation of ``diameter,'' define
\[
\delta(d,k)=\max\bigl\{\operatorname{diam}_{\mathrm{gr}}(P):P\subseteq[0,k]^d\text{ is a lattice polytope}\bigr\}.
\]
This definition includes polytopes of smaller affine dimension than \(d\). The supplied record does not explicitly specify the diameter convention or whether full dimension is required; these choices require source confirmation.

For lattice vectors \(v_1,\ldots,v_m\in\Z^d\), define the lattice zonotope
\[
Z(v_1,\ldots,v_m)=\sum_{i=1}^{m}[0,v_i]
=\left\{\sum_{i=1}^{m}\lambda_i v_i:0\leq\lambda_i\leq1\right\}.
\]
Here the sum is a Minkowski sum of line segments. This is the conventional zonotope interpretation of the source phrase ``Minkowski sum of lattice vectors''; that phrase does not explicitly mention segments and needs confirmation.

\textbf{Supplied conjectural questions, under these interpretations.} For every \(d,k\in\Z_{\geq1}\), do there exist \(m\in\Z_{\geq1}\), lattice vectors \(v_1,\ldots,v_m\in\Z^d\), and a translation vector \(t\in\Z^d\) such that
\[
Q=t+Z(v_1,\ldots,v_m)\subseteq[0,k]^d,
\qquad \operatorname{diam}_{\mathrm{gr}}(Q)=\delta(d,k)?
\]
In particular, the supplied record asks whether, for every \(d,k\in\Z_{\geq1}\) satisfying the strict inequality \(k<2d\),
\[
\delta(d,k)=\left\lfloor\frac{(k+1)d}{2}\right\rfloor.
\]
Both assertions are presented as questions, not established results. Their precise intended formulation remains subject to the qualifications above.
\subsection{Short English statement}
The supplied conjecture asks whether the largest graph diameter of a lattice polytope in an integer cube is attained by a translated lattice zonotope, and proposes an exact formula when \(k<2d\). Graph diameter, segment generators, and the inclusion of lower-dimensional polytopes are explicit interpretations requiring source confirmation.
\subsection{Sources}
\begin{itemize}
\item Supplied local UnsolvedMath registry record: permanent ID \(30003469\), title ``Maximum Diameters of Lattice Polytopes in Cubes,'' source problem code \texttt{catalogue entry}, in the Oberwolfach Reports Open Problems collection. The record gives proposed year 2017 and supplies no proposer attribution.
\item Supplied report citation: \emph{Discrete Geometry (2018)}, Oberwolfach Reports, DOI 10.4171/owr/2017/19, pp. 1135--1205. \url{https://doi.org/10.4171/owr/2017/19}. The supplied background describes these reports as documenting mathematical workshops and their open-problem discussions. The citation's 2018 label and the record's 2017 proposed year are retained without reconciliation.
\item Supplied literature reference: ``Deza and Pournin on diameter bounds.'' \url{https://arxiv.org/abs/1610.00341}.
\item Supplied literature reference: ``Chadder and Deza on exact small cases.'' \url{https://doi.org/10.1016/j.dam.2019.10.026}.
\item Literature assessment dated exactly \texttt{2026-08-21}: the supplied curation labels the problem \texttt{partially\_solved} and states that the conjectured exact formula remains unresolved, while strong upper bounds and computer-assisted exact values are known for several small parameter pairs. This is supplied dated curation derived from the OWR collection, not an independently verified literature assessment. The supplied background explicitly requests independent verification of time-sensitive status and cited claims.
\item The registry labels the statement \texttt{corrected\_verified}, explaining that mathematical content was preserved while notation, typography, and sentence structure were repaired for publication. This supplied editorial label is not independent mathematical verification of this draft.
\end{itemize}
\end{document}
